\documentclass{article}
\usepackage{graphicx} % Required for inserting images
\usepackage{amssymb}
\usepackage{amsmath}
\usepackage{mathtools}
\usepackage{url}

\newcommand{\numRacks}[0]{\ensuremath{\# r}}
\newcommand{\numTrailers}[0]{\ensuremath{\# t}}
\newcommand{\numJigs}[0]{\ensuremath{\# j}}
\newcommand{\numHangars}[0]{\ensuremath{\# h}}
\newcommand{\numBelugas}[0]{\ensuremath{\# b}}

\title{Beluga Domain in JANI}

\author{Chaahat Jain}

\date{2023}

\setlength{\parindent}{0pt}

\begin{document}

\maketitle

\tableofcontents

\section{Introduction}
This document provides an overview of JANI~\cite{budde:etal:tacas-17} model encodings of Beluga introduced by airbus for TUPLES.

%%%%%%%%%%%%%%%%%%%%%%%%%%%%%%%%%%%%%%%%%%%%%%%%%%%%%%%%%%%%%%%%%%%%%%%%%%%%%%%%%%%%%%%%%%%%%%%%%%%%%%%%%%%%%%%%%%%%%%%%
\section{Beluga}

Beluga is a TUPLES domain regarding transporting parts from a Beluga flight to a factory in a pre-determined order.
Since each jig can hold atmost one part, we consider jigs and parts to be the same object.

\subsection{Objects}
In this use case, we have the following objects:
\begin{itemize}
    \item Jigs: These are the parts that need to be delivered to the factory.
    \item Trailer: These are used to move jigs from one location to another.
    \item Beluga: The flight carrying the jigs to be delivered to the factory.
    \item Hangar: These are used to deliver the trailer carrying a jig to the factory.
    \item Storage: This is a location where multiple trailers start of at initially.
    \item Racks: These are used to store jigs. Jigs can be pushed/popped from either side of the rack. Only one trailer can stay on one side of the rack.
    \item Truck: A temporary stop for moving between two racks.
\end{itemize}

\subsection{State Variables}
A \textbf{state} in Beluga corresponds to a tuple of the variables described below.\\

\paragraph*{Jigs:} For each jig, we have three variables. $jig_j$ is used to denote the position of the jig. 
$front\_jig\_j$ is used to denote the jig in front of this jig. $back\_jig\_j$ is used to denote the jig behind this jig. The default value when no jig is present is $-1$.\\

\paragraph*{Trailers:} For each trailer, we have two variables. $trailer_t$ is used to denote the position of the trailer. 
$empty\_trailer_t$ is used to check if the trailer is empty ($1$ when trailer is empty). \\

\paragraph*{Beluga:} For each beluga flight, we have two variables. $beluga_b$ is used to denote the number of jigs on the beluga. 
$beluga_b\_phase$ is used to denote whether the beluga in in loading or unloading phase.\\

\paragraph*{Hangar:} For each hangar, we have one variable. $hangar_h$ is used to denote if the hangar is empty.\\

\paragraph*{Racks:} For each rack, we have four variables. $rack\_load_r$ is used to denote the remaining capacity of the rack. $num\_jigs\_on\_rack_r$ denotes how many jigs are on a rack. 
$front\_rack_r$ and $back\_rack_r$ are used to find the first and last jigs on a rack.\\

\paragraph*{Truck:} We have one variable for the truck, denoting whether the truck is currently free.\\

\paragraph*{Factory}: We have one variable for how many jigs have been delivered to the factory.

\subsection{Global Constants}
For any given state, we know the total number of racks, jigs, hangars, belugas and trailers. This is represented as $\numRacks, \numJigs, \numHangars, \numBelugas, \numTrailers$ respectively.

\subsection{Assignments}
For the position of trailer, the position of jigs and the order of sending jigs to the factory, we define the following assignments.\\

\paragraph*{Sending Jigs to Factory}: Our naming of jigs is according to the order in which they need to go to the factory. For example, $jig_0$ is the first jig required by the factory.\\

\paragraph*{Position of Jigs:} Every jig can either be on a rack, on the beluga or in a trailer. We interpret the beluga to be a stack. Positions of jigs on the beluga are encoded as distance from the bottom of the stack. As such, we can use the following function:
\begin{equation*}
    jig_j =\begin{cases*}
        \numJigs*b + i, & \text{jig on beluga b, at distance} i ~\text{from bottom}\\
        \numJigs*\numBelugas + r, & \text{jig on rack} {r} \\
        \numJigs*\numBelugas + \numRacks + t, & \text{jig on trailer} {t}
        \end{cases*}
\end{equation*}

\paragraph*{Position of Trailers:} Every trailer can either be in front of or behind a rack, or it can be on a hangar, or in storage, or in front of a beluga or in front of a truck.
\begin{equation*}
    trailer_t =\begin{cases*}
        -r - 1, & \text{trailer behind rack} {r}\\
        r, & \text{trailer in front of rack} {r} \\
        \numRacks + h, & \text{trailer on hangar} {h} \\
        \numRacks + \numHangars + b, & \text{trailer in front of beluga} {b}\\
        \numHangars + \numRacks + \numBelugas, & \text{trailer in storage}\\
        \numHangars + \numRacks + \numBelugas + 1, & \text{trailer in front of truck}
        \end{cases*}
\end{equation*}


\subsection{Actions}
The actions have been encoded exactly as presented in the Beluga Use Case document from: \url{https://drive.google.com/file/d/1ZsEBH6jiQV85ZAttSmqpsBO_R-l0bKeC/view}.

\subsection{Goal Condition and Safety}
The current \textbf{goal condition} we look at is to deliver all parts to the factory and then load the jigs back onto the beluga.
The \textbf{safety condition} currently checks that neither rack is filled more than $80\%$.

%%%%%%%%%%%%%%%%%%%%%%%%%%%%%%%%%%%%%%%%%%%%%%%%%%%%%%%%%%%%%%%%%%%%%%%%%%%%%%%%%%%%%%%%%%%%%%%%%%%%%%%%%%%%%%%%%%%%%%%%
\section{Beluga PDDL To JANI}
A simpler version of beluga is used for the PDDL version when training ASNets. In this section, we describe an equivalent encoding in JANI. 

\subsection{Changes to original version}
The following simplification are made:
\begin{itemize}
    \item Hangars have been totally removed. We have exactly 1 trailer. The trailer position does not matter since this trailer can always unload belugas and push jigs on racks.

    \item We do not load jigs back into beluga. We are only unloading from beluga and sending to factory.
    
    \item We can always unload from beluga as long as a trailer is empty. We can load onto a rack without worrying about this trailer position.
    
    \item We can take jig from rack and directly put on truck in one action. 
    
    \item Only items on trucks are sent to the factory.
\end{itemize}

\subsection{Action}
We consider the following actions:
\begin{enumerate}
    \item Unloading from Beluga 
    \begin{itemize}
        \item Removes the top jig of the Beluga stack as long as the trailer is empty.
    \end{itemize}
    \item Pushing to front of rack r 
    \begin{itemize}
        \item Only done when there is enough space on the racks to hold the jig.
        \item We split into two cases based on whether the rack is empty. 
        \item Variable for back of jig must be updated. Variables for front of rack must be updated. Trailer is now empty.
    \end{itemize}
    \item Pushing to back of rack r
    \begin{itemize}
        \item Only done when there is enough space on rack to hold the jig.
        \item We split into two cases based on whether the rack is empty. 
        \item Variable for front of jig must be updated. Variables for back of rack must be updated. Trailer is now empty.
    \end{itemize}
    \item Popping front of rack r to truck
    \begin{itemize}
        \item Only done when truck is empty.
        \item We split into two cases based on whether the rack will be empty after the pop. 
        \item Variable for front of rack must be updated. Truck is now not empty.
    \end{itemize}
    \item Popping back of rack r to truck     
    \begin{itemize}
        \item Only done when truck is empty.
        \item We split into two cases based on whether the rack will be empty after the pop. 
        \item Variable for back of rack must be updated. Truck is now not empty.
    \end{itemize}
    \item Pushing truck to front of rack r    
    \begin{itemize}
        \item Only done when truck is not empty and rack has enough space to hold jig.
        \item We split into two cases based on whether the rack is empty. 
        \item Variable for front of rack must be updated. Variable for back of jig must be updated. Truck is now empty.
    \end{itemize}
    \item Pushing truck to back of rack r   
    \begin{itemize}
        \item Only done when truck is not empty and rack has enough space to hold jig.
        \item We split into two cases based on whether the rack is empty. 
        \item Variable for back of rack must be updated. Variable for front of jig must be updated. Truck is now empty.
    \end{itemize}
    \item Sending to factory    
    \begin{itemize}
        \item Only done when truck is not empty and jig held in truck is next to be sent to factory.
        \item Truck is now empty. The number of parts delivered to factory increase by 1. We do not update jig position yet because this is simplification.
    \end{itemize}
\end{enumerate}

\bibliographystyle{plain}
\bibliography{bib}

\end{document}
